\chapter*{Abstract}
\addcontentsline{toc}{chapter}{Abstract}

Graph generation aims to learn distributions capable of producing new graphs
that preserve the structural properties of observed data. This thesis introduces
OGAD (On Graphons And Diffusion), a graph generative pipeline that uses
graphon-derived representations as an intermediate domain between finite graphs
and discrete diffusion. Starting from training graphs, OGAD estimates graphons
with Scalable Implicit Graphon Learning (SIGL), evaluates them on finite grids,
quantizes the resulting probability matrices into categorical tensors, and
models their distribution with a Discrete Denoising Diffusion Probabilistic
Model (D3PM). Generated tensors are converted back into probability grids from
which finite graphs are sampled.

OGAD is evaluated on five synthetic graph families and compared with DiGress
and classical probabilistic generators using Maximum Mean Discrepancy over
degree distributions, clustering coefficients, and normalized-Laplacian
spectra. Results are strongly descriptor- and dataset-dependent. OGAD is
particularly effective at preserving clustering structure, achieving the
lowest clustering discrepancy on Stochastic Block Model and expected-degree
graphs, while degree and spectral properties are often better captured by
direct graph diffusion or specialized baselines. Overall, graphon-mediated
diffusion emerges as a viable alternative for graph generation with a distinct
structural bias.